\chapter{Diagonalization}

\section{Similar matrices}

\begin{definition}
$n\times n$ matrices $A$ and $B$ are \emph{similar} if there is an invertible $P$ with $P^{-1}AP=B$.
\end{definition}

Similarity is an equivalence relation (reflexive with $P=I$; symmetric with $P^{-1}$; transitive since
$(PQ)^{-1}A(PQ)=Q^{-1}(P^{-1}AP)Q$). In Chapter 12 we will see that similar matrices represent the same
linear operator in different bases, which explains why they share so much.

\begin{theorem}
If $A$ and $B$ are similar, then they have the same characteristic polynomial, hence the same eigenvalues
with the same algebraic multiplicities, and the same determinant, trace and rank.
\end{theorem}
\begin{proof}
$\det(B-\lambda I)=\det(P^{-1}AP-\lambda P^{-1}P)=\det\big(P^{-1}(A-\lambda I)P\big)
=\det(P^{-1})\det(A-\lambda I)\det P=\det(A-\lambda I)$. Determinant and trace are (up to sign) coefficients of
the characteristic polynomial. For rank, multiplying by invertible matrices does not change the null space
dimension: $B\vec x=\vec 0\iff A(P\vec x)=\vec0$, so $\Null B=P^{-1}\Null A$.
\end{proof}

\begin{important}[Pitfall]
The converse fails: $I$ and $\begin{bmatrix}1&1\\0&1\end{bmatrix}$ have the same characteristic polynomial
$(1-\lambda)^2$ but are not similar (the only matrix similar to $I$ is $P^{-1}IP=I$).
\end{important}

\section{The Diagonalization Theorem}

\begin{definition}
$A$ is \emph{diagonalizable} if it is similar to a diagonal matrix $D$: $P^{-1}AP=D$. We say $P$
\emph{diagonalizes} $A$.
\end{definition}

\begin{theorem}[Diagonalization Theorem]
An $n\times n$ matrix $A$ is diagonalizable (over $\R$) if and only if there is a basis of $\R^n$ consisting
of eigenvectors of $A$. In that case, if $\{\vec v_1,\dots,\vec v_n\}$ is such a basis with
$A\vec v_i=\lambda_i\vec v_i$, then
\[
  P=[\vec v_1\ \cdots\ \vec v_n],\qquad P^{-1}AP=D=\diag(\lambda_1,\dots,\lambda_n).
\]
\end{theorem}
\begin{proof}
$P^{-1}AP=D$ iff $AP=PD$ (with $P$ invertible). Column-by-column, $AP=[A\vec v_1\cdots A\vec v_n]$ and
$PD=[\lambda_1\vec v_1\cdots\lambda_n\vec v_n]$. So $AP=PD$ iff each $\vec v_i$ satisfies
$A\vec v_i=\lambda_i\vec v_i$; and $P$ is invertible iff its columns form a basis of $\R^n$ (Invertible Matrix
Theorem). Since basis vectors are non-zero, they are eigenvectors.
\end{proof}

\begin{theorem}[Multiplicity test]
$A$ is diagonalizable over $\R$ if and only if all roots of $C(\lambda)$ are real and $g_\lambda=a_\lambda$ for
every eigenvalue $\lambda$. In particular, an $n\times n$ matrix with $n$ distinct real eigenvalues is
diagonalizable.
\end{theorem}
\begin{proof}
Combining bases of the eigenspaces gives a linearly independent set (Chapter 9) with $\sum g_\lambda$ vectors.
Since $g_\lambda\le a_\lambda$ and $\sum a_\lambda\le n$ (with equality iff all roots are real), we get $n$ vectors
exactly when all roots are real and $g_\lambda=a_\lambda$ for all $\lambda$. Conversely, any eigenbasis contains at
most $g_\lambda$ vectors from each $E_\lambda$.
\end{proof}

\begin{important}[Algorithm]
\begin{enumerate}
  \item Find and factor $C(\lambda)$. If some root is non-real, $A$ is not diagonalizable over $\R$.
  \item For each eigenvalue find a basis of $E_\lambda$. If some $g_\lambda<a_\lambda$, stop: not diagonalizable.
  \item Put the eigenvectors as the columns of $P$ and the corresponding eigenvalues, \emph{in the same order},
        on the diagonal of $D$.
\end{enumerate}
\end{important}

\begin{example}
For $A=\begin{bmatrix}4&1\\2&3\end{bmatrix}$ (Chapter 9) we found $E_2=\Span\{(1,-2)\}$ and $E_5=\Span\{(1,1)\}$.
So
\[
  P=\begin{bmatrix}1&1\\-2&1\end{bmatrix},\quad D=\begin{bmatrix}2&0\\0&5\end{bmatrix},\quad
  P^{-1}=\frac13\begin{bmatrix}1&-1\\2&1\end{bmatrix}.
\]
Check $AP=\begin{bmatrix}4-2&4+1\\2-6&2+3\end{bmatrix}=\begin{bmatrix}2&5\\-4&5\end{bmatrix}$ and
$PD=\begin{bmatrix}2&5\\-4&5\end{bmatrix}$. Equal, so $P^{-1}AP=D$.
\end{example}

In the eigenvector basis the action of $A$ is pure stretching: the parallelogram on
$\vec v_1 = (1,-2)$ and $\vec v_2 = (1,1)$ is sent to the one on $2\vec v_1$ and $5\vec v_2$. That is what
$D = \diag(2,5)$ says.
\begin{center}
\begin{tikzpicture}[scale=0.42,>={Stealth}]
  \draw[->,gray] (-1,0) -- (8,0); \draw[->,gray] (0,-4.6) -- (0,5.6);
  \fill[orange!20] (0,0) -- (2,-4) -- (7,1) -- (5,5) -- cycle;
  \draw[orange!80!black] (0,0) -- (2,-4) -- (7,1) -- (5,5) -- cycle;
  \fill[blue!20] (0,0) -- (1,-2) -- (2,-1) -- (1,1) -- cycle;
  \draw[blue!70!black] (0,0) -- (1,-2) -- (2,-1) -- (1,1) -- cycle;
  \draw[->,thick,blue!70!black] (0,0) -- (1,-2) node[left] {$\vec v_1$};
  \draw[->,thick,blue!70!black] (0,0) -- (1,1) node[above left] {$\vec v_2$};
  \draw[->,thick,orange!80!black] (0,0) -- (2,-4) node[below] {$A\vec v_1 = 2\vec v_1$};
  \draw[->,thick,orange!80!black] (0,0) -- (5,5) node[above] {$A\vec v_2 = 5\vec v_2$};
\end{tikzpicture}
\end{center}

\begin{example}
For $A=\begin{bmatrix}0&1&1\\1&0&1\\1&1&0\end{bmatrix}$ we found $a_{-1}=g_{-1}=2$ and $a_2=g_2=1$, so $A$ is
diagonalizable with
\[
  P=\begin{bmatrix}-1&-1&1\\1&0&1\\0&1&1\end{bmatrix},\qquad D=\begin{bmatrix}-1&0&0\\0&-1&0\\0&0&2\end{bmatrix}.
\]
(Check one column: $A(-1,0,1)=(1,0,-1)=-1\cdot(-1,0,1)$.)
\end{example}

\begin{example}[Not diagonalizable]
$A=\begin{bmatrix}3&1&0\\0&3&0\\0&0&5\end{bmatrix}$ is upper triangular, so its eigenvalues are $3$ ($a_3=2$) and
$5$ ($a_5=1$). $A-3I=\begin{bmatrix}0&1&0\\0&0&0\\0&0&2\end{bmatrix}$ has rank $2$, so
$g_3=3-2=1<2$. Not diagonalizable.
\end{example}

\section{Powers of matrices}

If $A=PDP^{-1}$, then $A^k=(PDP^{-1})(PDP^{-1})\cdots(PDP^{-1})=PD^kP^{-1}$ because the inner
$P^{-1}P$ factors cancel; and $D^k$ is computed entrywise on the diagonal.

\begin{example}
For $A=\begin{bmatrix}4&1\\2&3\end{bmatrix}$:
\[
A^k=PD^kP^{-1}=\begin{bmatrix}1&1\\-2&1\end{bmatrix}\begin{bmatrix}2^k&0\\0&5^k\end{bmatrix}\frac13\begin{bmatrix}1&-1\\2&1\end{bmatrix}
=\frac13\begin{bmatrix}2^k&5^k\\-2^{k+1}&5^k\end{bmatrix}\begin{bmatrix}1&-1\\2&1\end{bmatrix}
\]
\[
=\frac13\begin{bmatrix}2^k+2\cdot5^k & 5^k-2^k\\ 2\cdot5^k-2^{k+1} & 5^k+2^{k+1}\end{bmatrix}.
\]
Check $k=1$: $\frac13\begin{bmatrix}2+10&5-2\\10-4&5+4\end{bmatrix}=\begin{bmatrix}4&1\\2&3\end{bmatrix}$. Check $k=0$:
$\frac13\begin{bmatrix}3&0\\0&3\end{bmatrix}=I$.
\end{example}

\begin{example}[A Markov chain]
Each year $10\%$ of city residents move to the suburbs and $20\%$ of suburban residents move to the city.
If $\vec x_k=(c_k,s_k)$ are the population fractions in year $k$, then $\vec x_{k+1}=T\vec x_k$ with
$T=\begin{bmatrix}0.9&0.2\\0.1&0.8\end{bmatrix}$.
\begin{center}
\begin{tikzpicture}[>={Stealth},node distance=4cm]
  \node[draw,circle,minimum size=1.5cm,fill=blue!10] (c) {city};
  \node[draw,circle,minimum size=1.5cm,fill=green!10,right of=c] (s) {suburbs};
  \draw[->,thick] (c) to[bend left=25] node[above] {$0.1$} (s);
  \draw[->,thick] (s) to[bend left=25] node[below] {$0.2$} (c);
  \draw[->,thick] (c) to[loop left] node[left] {$0.9$} (c);
  \draw[->,thick] (s) to[loop right] node[right] {$0.8$} (s);
\end{tikzpicture}
\end{center}
Column $j$ of $T$ lists where the people currently in place $j$ go. Its characteristic polynomial is
$(0.9-\lambda)(0.8-\lambda)-0.02=\lambda^2-1.7\lambda+0.7=(\lambda-1)(\lambda-0.7)$.
Eigenvectors: $T-I=\begin{bmatrix}-0.1&0.2\\0.1&-0.2\end{bmatrix}$ gives $(2,1)$; $T-0.7I=\begin{bmatrix}0.2&0.2\\0.1&0.1\end{bmatrix}$
gives $(1,-1)$. Writing $\vec x_0=a(2,1)+b(1,-1)$,
\[
  \vec x_k=T^k\vec x_0=a(2,1)+b(0.7)^k(1,-1)\longrightarrow a(2,1).
\]
Since $c_k+s_k=1$ always, and the $(1,-1)$ term contributes $0$ to the sum, $3a=1$; the long-run state is
$\left(\frac23,\frac13\right)$ regardless of the starting distribution.
\end{example}

\begin{aside}
Diagonalization is also used to solve linear recurrences: the Fibonacci numbers satisfy
$\begin{bmatrix}F_{k+1}\\F_k\end{bmatrix}=\begin{bmatrix}1&1\\1&0\end{bmatrix}\begin{bmatrix}F_k\\F_{k-1}\end{bmatrix}$, and diagonalizing
this matrix (eigenvalues $\frac{1\pm\sqrt5}{2}$) gives Binet's closed formula.
\end{aside}

\section*{Summary}
\begin{itemize}
  \item Similar matrices ($B=P^{-1}AP$) share characteristic polynomial, eigenvalues, determinant, trace, rank.
  \item $A$ is diagonalizable iff $\R^n$ has a basis of eigenvectors iff all eigenvalues are real with
        $g_\lambda=a_\lambda$. Distinct eigenvalues suffice.
  \item $P=[\text{eigenvectors}]$, $D=\diag(\text{eigenvalues in the same order})$, $AP=PD$.
  \item $A^k=PD^kP^{-1}$.
\end{itemize}
